\documentclass[11pt]{article}
\usepackage[T1]{fontenc}
\usepackage{lmodern}
\usepackage[margin=1.05in]{geometry}
\usepackage{amsmath,amssymb,amsthm,mathtools}
\usepackage{booktabs,array,microtype}
\usepackage[hidelinks]{hyperref}
\ifdefined\pdfinfoomitdate\pdfinfoomitdate=1\fi
\ifdefined\pdfsuppressptexinfo\pdfsuppressptexinfo=15\fi
\ifdefined\pdftrailerid\pdftrailerid{}\fi
\hypersetup{pdftitle={Principal components and singularities of orthogonal frame varieties},
pdfauthor={Ying Xie},pdfsubject={Commutative algebra and algebraic geometry},
pdfkeywords={orthogonal frames, defining ideals, canonical modules, linkage, rational singularities},
pdfcreator={},pdfproducer={},pdfstartview={}}
\newtheorem{theorem}{Theorem}[section]
\newtheorem{proposition}[theorem]{Proposition}
\newtheorem{lemma}[theorem]{Lemma}
\newtheorem{corollary}[theorem]{Corollary}
\theoremstyle{definition}
\newtheorem{definition}[theorem]{Definition}
\newtheorem{example}[theorem]{Example}
\theoremstyle{remark}
\newtheorem{remark}[theorem]{Remark}
\numberwithin{equation}{section}
\DeclareMathOperator{\Mat}{Mat}
\DeclareMathOperator{\rank}{rank}
\DeclareMathOperator{\diag}{diag}
\DeclareMathOperator{\Sym}{Sym}
\DeclareMathOperator{\Hom}{Hom}
\DeclareMathOperator{\Ext}{Ext}
\DeclareMathOperator{\Spec}{Spec}
\DeclareMathOperator{\Proj}{Proj}
\DeclareMathOperator{\Gr}{Gr}
\DeclareMathOperator{\OGr}{OGr}
\DeclareMathOperator{\GL}{GL}
\DeclareMathOperator{\SL}{SL}
\DeclareMathOperator{\SO}{SO}
\DeclareMathOperator{\Tot}{Tot}
\DeclareMathOperator{\depth}{depth}
\DeclareMathOperator{\pd}{pd}
\DeclareMathOperator{\reg}{reg}
\DeclareMathOperator{\indeg}{indeg}
\DeclareMathOperator{\codim}{codim}
\newcommand{\kk}{K}
\newcommand{\CC}{\mathbb C}
\newcommand{\RR}{\mathbb R}
\newcommand{\PP}{\mathbb P}
\newcommand{\OO}{\mathcal O}
\newcommand{\UU}{\mathcal U}
\newcommand{\XX}{\mathsf X}
\newcommand{\II}{\mathsf I}
\newcommand{\Dpr}{D_{\mathrm{prime}}}
\newcommand{\DCI}{D_{\mathrm{CI}}}
\newcommand{\DUFD}{D_{\mathrm{UFD}}}
\newcommand{\T}{\mathsf T}
\newcommand{\Schur}{\mathbf S}
\newcommand{\doi}[1]{\href{https://doi.org/#1}{\nolinkurl{#1}}}
\setlength{\emergencystretch}{1.5em}
\title{Principal components and singularities\\of orthogonal frame varieties}
\author{Ying Xie\\[3pt]\small Kennesaw State University\\
\small\texttt{yxie2@kennesaw.edu}}
\date{}

\begin{document}
\maketitle

\begin{abstract}
We study the defining ideals and singularities of varieties of pairwise
orthogonal vectors. For $q=\binom{k}{2}$ and $h=q+k$, we determine the
ideal of the principal component of the variety of $h$ orthogonal vectors
in dimension $2q$ over $\CC$. Besides the Gram quadrics, the ideal has
generators in the single degree $q(k+1)$, constructed by Gale duality
from a Veronese determinant and indexed by two orthogonal-group
representations. In the first case, six vectors in dimension six, this
gives fifteen quadrics and $330$ explicit equations of degree twelve.
In every characteristic different from two, the principal components
in this family are normal, satisfy $S_k$ but not $S_{k+1}$, and have
depth one below their dimension. We determine their deficiency modules
and regularity. We also construct normal, and frequently factorial,
orthogonal frame varieties at the prime threshold that fail to have
rational singularities in characteristic zero and fail to be $F$-rational
in every odd characteristic. The corresponding numbers of columns form
a set of natural density one half. An exact local chart identifies the
generic transverse singularities with cones over smooth complete
intersections of quadrics and gives their classification. These results
answer the first defining-ideal case and the rationality question posed
by Casabella and Sammartano, and give normal but non-Cohen--Macaulay
principal components.
\end{abstract}

\section{Introduction}

Let $E$ be a nondegenerate quadratic space of dimension $M$ over an
algebraically closed field $K$ of characteristic different from two.
For a generic matrix $V=[v_1\ \cdots\ v_h]\in\Mat_{M,h}$, set
\[
 S=K[v_{ai}],\qquad
 \II(M,h)=(\langle v_i,v_j\rangle:1\leq i<j\leq h),\qquad
 \XX(M,h)=\Spec S/\II(M,h).
\]
When $M\geq h$, the closure of the locus on which all column norms are
nonzero is an irreducible component. We denote it by $C(M,h)$ and call
it the \emph{principal component}; its prime ideal is $P(M,h)$.

Casabella and Sammartano determine the irreducible components of
orthogonal frame varieties and the thresholds for complete intersection,
primality, and normality \cite{CS}. Their Questions 8.1, 8.3, and 8.4 ask,
respectively, about rational singularities at the prime threshold,
defining equations of the principal component, and its normality and
Cohen--Macaulayness. The first unresolved defining-ideal case identified
there is $(M,h)=(6,6)$.

Our results concern two families adjacent to these thresholds. The first
lies at a complete-intersection threshold where the full scheme has
three components. The second lies at a prime threshold. Both are governed
by the matrix of evaluations of the monomials $z_i z_j$, $i<j$, at a
collection of points. In the first family its determinant describes the
boundary of the principal component; in the second its kernel gives
the transverse quadrics.

\subsection{The principal component}

For the first family fix
\begin{equation}\label{eq:triangular}
 k\geq3,\qquad q=\binom{k}{2},\qquad
 h=q+k=\binom{k+1}{2},\qquad M=2q,
 \qquad d_0=q(k+1)=h(k-1).
\end{equation}
The symbol $d_0$ denotes a degree, not an ambient dimension. Put
\begin{equation}\label{eq:generator-number}
 b_k=\prod_{1\leq i<j\leq q}
       \frac{2(k+1)+2q-i-j}{2q-i-j}.
\end{equation}

\begin{theorem}\label{thm:ideal-intro}
Work over $\CC$, with the parameters in \eqref{eq:triangular}.
The module $P(2q,h)/\II(2q,h)$ is generated by its degree-$d_0$ part,
which has dimension $2b_k$, and it has no nonzero element of smaller
degree. Its generators are obtained as follows. Let $G_k$ be the
Gale-dual Veronese polynomial of Definition~\ref{def:Gale}. As $U$ ranges
over the maximal-isotropic $q$-dimensional row spaces in $K^{2q}$, the
polynomials $G_k(UV)$ span two spaces of dimension $b_k$, one for each
component of $\OGr(q,2q)$. Their images are independent modulo
$\II(2q,h)$ and minimally generate $P(2q,h)/\II(2q,h)$.
\end{theorem}

The construction is finite and explicit: one takes coefficients on a
dense skew-matrix chart in each orthogonal Grassmannian and extracts
a basis. The first numerical cases are
\[
\begin{array}{c|r|r|r|r}
 k&M&h&d_0&2b_k\\\hline
 3&6&6&12&330\\
 4&12&10&30&163324304\\
 5&20&15&60&9627912669442441500
\end{array}
\]
For $(6,6)$, Section~\ref{sec:six} gives all $330$ equations directly
as the coefficients of two determinants, without a basis-extraction step.

For the homological statements write
\begin{equation}\label{eq:homological-notation}
 L=2qh,\qquad e=\binom h2,\qquad N=L-e,
 \qquad a_R=2e-L.
\end{equation}
All polynomial variables have degree one, and $T(s)_j=T_{j+s}$.

\begin{theorem}\label{thm:depth-intro}
In every characteristic different from two, $C=C(2q,h)$ is normal
and satisfies $S_k$ but not $S_{k+1}$. Its non-Cohen--Macaulay locus is
\[
 Z=\{V:V^{\mathsf t}BV=0,\ \rank V\leq q-1\},
\]
where $B$ is a matrix of the quadratic form. This locus has codimension
$k+1$ in $C$, and
\[
 \depth\OO_{C,x}=\dim\OO_{C,x}-1\qquad(x\in Z).
\]
Let $Q_+$ and $Q_-$ be the primes of the two maximal-isotropic
components, set $J=Q_+\cap Q_-$, and put $A=S/P$ and
$Z_0=S/(Q_++Q_-)$. Then
\begin{align*}
 \Ext^e_S(A,S(-L))&=(J/\II)(a_R),\\
 \Ext^{e+1}_S(A,S(-L))&=Z_0(a_R),
\end{align*}
and all other Ext modules vanish. Consequently
\[
 \depth A=N-1,\qquad \pd_S A=e+1,\qquad \reg A=e-1.
\]
The last nonzero module in its minimal graded free resolution is
$S(-2e)$.
\end{theorem}

Thus the principal components in this family are normal but do not have
rational singularities in characteristic zero and are not $F$-rational
in odd characteristic: both properties imply Cohen--Macaulayness.
The defining-ideal theorem is restricted to characteristic zero because
its proof uses semisimplicity of the relevant representations. The
depth and normality assertions have no such restriction.

\subsection{Counterexamples at the prime threshold}

The prime threshold of \cite{CS} is
\begin{equation}\label{eq:prime-threshold}
 \Dpr(h)=\min\left\{
 2\left\lfloor\frac{2h+1-\sqrt{8h+1}}2\right\rfloor+2,
 2\left\lfloor\frac{2h+1-\sqrt{8h-7}}2\right\rfloor+1
 \right\}.
\end{equation}

\begin{theorem}\label{thm:negative-intro}
For $k\geq4$, let
\begin{equation}\label{eq:bad-range}
 \binom k2+2\leq h\leq\left\lfloor\frac{k^2}{2}\right\rfloor,
 \qquad M=2h-2k+1.
\end{equation}
Then $M=\Dpr(h)$ and $\XX(M,h)$ is a normal integral complete
intersection. It does not have rational singularities over $\CC$ and
is not $F$-rational in any odd characteristic. If $h<k^2/2$, it is
not log canonical over $\CC$ and is not $F$-pure in any odd
characteristic. The set of integers $h$ in \eqref{eq:bad-range} has
natural density $1/2$. The same density is obtained by imposing
$h-\binom k2\geq4$, in which case the coordinate rings are factorial.
\end{theorem}

This gives a negative answer to \cite[Question 8.1]{CS}. The first
example is $\XX(9,8)$. At a generic maximal-isotropic frame it has a
transverse cone over a smooth elliptic quartic. The next example,
$\XX(15,12)$, is already not log canonical. The generic transverse
classification, including the positive cases, is given in
Theorem~\ref{thm:transverse}. Those positive assertions concern the
generic point of one stratum, rather than all points of the frame
variety.

Theorem~\ref{thm:ideal-intro} answers \cite[Question 8.3]{CS} for the
entire family \eqref{eq:triangular}, including the stated first case.
Theorem~\ref{thm:depth-intro} gives negative answers to the universal
Cohen--Macaulay and rationality assertions in Question 8.4 and proves
normality for the same family. It does not classify principal components
for arbitrary $(M,h)$.

\section{Components, canonical modules, and a local chart}\label{sec:prelim}

\subsection{The three-component decomposition}

For the parameters \eqref{eq:triangular}, the component and
complete-intersection theorems of Casabella--Sammartano
\cite[Theorems 1 and 5.8]{CS} give
\begin{equation}\label{eq:three-components}
 \II=P\cap Q_+\cap Q_-,\qquad
 \XX(2q,h)=C\cup Y_+\cup Y_-.
\end{equation}
Here $Y_\pm$ are the closures of frames spanning maximal isotropic
$q$-planes of the two types. The ideal $\II$ is a radical complete
intersection of $e$ quadrics; all three components have dimension
\begin{equation}\label{eq:dimension-N}
 N=qh+\frac{q(q-1)}2=L-e.
\end{equation}
Indeed, substitution of $8h+1=(2k+1)^2$ into the threshold formulas
gives $\DCI(h)=2q$ and $\Dpr(h)=2q+1$.

We use the following additional input from
Hochster--Jeffries--Pandey--Singh \cite[Theorem 7.13]{HJPS}:
the rings $S/Q_+$, $S/Q_-$, and $Z_0=S/(Q_++Q_-)$ are
Cohen--Macaulay domains, and
\begin{equation}\label{eq:intersection}
 Q_++Q_-=(V^{\mathsf t}BV)+I_q(V),\qquad
 \dim Z_0=N-k-1.
\end{equation}
The parentheses in \eqref{eq:intersection} denote all Gram entries,
including the diagonal; $I_q(V)$ denotes the ideal of $q$-minors.
Thus $\Spec Z_0$ is the rank-at-most-$(q-1)$ isotropic locus.

Canonical modules and graded local duality are used with the conventions
of \cite{BH}. In particular, for an equidimensional quotient $D$ of
$S$ of codimension $e$, we write
$\omega_D=\Ext^e_S(D,S(-L))$, without requiring $D$ to be
Cohen--Macaulay. With $R=S/\II$, one has $\omega_R=R(a_R)$.

\begin{lemma}\label{lem:linkage-canonical}
Let $J=Q_+\cap Q_-$. There are graded isomorphisms
\begin{equation}\label{eq:canonical-linkage}
 \omega_{S/J}=\omega_{S/Q_+}\oplus\omega_{S/Q_-},\qquad
 P/\II=\omega_{S/J}(-a_R).
\end{equation}
In particular $P/\II$ is a Cohen--Macaulay $S$-module of dimension $N$.
\end{lemma}

\begin{proof}
Apply $\Ext_S^\bullet(-,S(-L))$ to
\begin{equation}\label{eq:union}
 0\longrightarrow S/J\longrightarrow
 S/Q_+\oplus S/Q_-\longrightarrow Z_0\longrightarrow0.
\end{equation}
The codimension of $Z_0$ in $S$ is $e+k+1$, so its Ext modules
vanish in degrees $e$ and $e+1$. This proves the first isomorphism.
The radical decomposition gives $\II:J=P$. Change of rings for the
complete intersection yields
\[
 \omega_{S/J}=\Hom_R(S/J,\omega_R)=(P/\II)(a_R).
\]
This is the canonical-module form of complete-intersection linkage
\cite{PS}. Each summand in the first isomorphism is Cohen--Macaulay
of dimension $N$, proving the last assertion.
\end{proof}

\subsection{An exact chart at an isotropic frame}

The next chart will be used in both families. Write a split quadratic
space as $K^q\oplus K^q\oplus K^m$ with pairing
\[
 \langle(u,v,z),(u',v',z')\rangle
 =u^{\mathsf t}v'+v^{\mathsf t}u'+z^{\mathsf t}z'.
\]
For $m=0$ or $1$, the rank-$q$ isotropic stratum is maximal isotropic.

\begin{lemma}\label{lem:chart}
Let $h=q+k$ and $M=2q+m$. On a chart meeting the rank-$q$ isotropic
stratum, the orthogonality equations are, after an invertible change
of coordinates and elimination of variables,
\begin{equation}\label{eq:chart}
 \sum_{a=1}^m z_{ai}z_{aj}
 -\sum_{r=1}^q\lambda_r y_{ri}y_{rj}=0
 \qquad(1\leq i<j\leq k).
\end{equation}
The coordinates $Y=(y_{ri})$ are unrestricted. In addition there are
$q^2+\binom q2+mq$ free coordinates, with a $q\times q$ determinant
inverted. The isotropic stratum in this chart is $\lambda=z=0$.
\end{lemma}

\begin{proof}
Choose $q$ independent columns whose projection onto the first summand
is invertible, and denote this top block by $D$. The orthogonal change
$\diag(D^{-1},D^{\mathsf t},I_m)$ makes it the identity and retains
the entries of $D$ as independent chart coordinates. Write the columns
in blocks
\[
 A=\begin{bmatrix}I_q\\H\\C\end{bmatrix},\qquad
 B'=\begin{bmatrix}Y\\T\\W\end{bmatrix}.
\]
The first set of equations is
$H+H^{\mathsf t}+C^{\mathsf t}C=\diag(\lambda)$.
It eliminates the symmetric part of $H$. The equations
$A^{\mathsf t}BB'=0$ for the ambient form, equivalently the pairings
between columns of $A$ and $B'$, give
$T=-H^{\mathsf t}Y-C^{\mathsf t}W$.
With $z=W-CY$, direct substitution in the Gram matrix of $B'$ gives
\[
 (B')^{\mathsf t}BB'=z^{\mathsf t}z-Y^{\mathsf t}\diag(\lambda)Y.
\]
Its off-diagonal entries are \eqref{eq:chart}. The changes used are
invertible on the chosen chart, so this describes the local scheme,
including its transverse parameters. If $\lambda=0$, the columns of
$A$ are isotropic; $z=0$ says that the remaining columns belong to
their span. This proves the last assertion.
\end{proof}

For a $q\times k$ matrix $Y$, define its hollow-quadratic evaluation
matrix by
\begin{equation}\label{eq:E}
 E(Y)_{(ij),r}=y_{ri}y_{rj},\qquad i<j.
\end{equation}
It has $\binom k2$ rows and $q$ columns. The case $m=0$ of
\eqref{eq:chart} is $E(Y)\lambda=0$; when $m=1$, it is
$E(Y)\lambda=(z_i z_j)_{i<j}$.

\section{The defining ideal in the triangular family}\label{sec:ideal}

Throughout this section the field is $\CC$ and the parameters are
\eqref{eq:triangular}. Use the split form
$B=\left[\begin{smallmatrix}0&I_q\\I_q&0\end{smallmatrix}\right]$.

\subsection{The Gale-dual Veronese polynomial}

Let $\nu_2$ list the $h=\binom{k+1}{2}$ quadratic monomials in $k$
variables, in a fixed order. For an $h\times k$ matrix $K_0$ with
rows $\kappa_i$, put
\begin{equation}\label{eq:Psi}
 \Psi_k(K_0)=\det[\nu_2(\kappa_i)]_{i=1}^h.
\end{equation}
The vectors in the determinant may be taken as columns; an overall
sign from the ordering will be immaterial. Since
$\det\Sym^2(g)=\det(g)^{k+1}$,
\begin{equation}\label{eq:Psi-character}
 \Psi_k(K_0g)=\det(g)^{k+1}\Psi_k(K_0).
\end{equation}
The first fundamental theorem for $\SL_k$ expresses this polynomial
in the maximal minors of $K_0$; see the invariant-theoretic discussion
in \cite{HJPS}. Its degree in these minors is $k+1$.

\begin{definition}\label{def:Gale}
Apply the complementary-minor identification between the Pl\"ucker
coordinate rings of $\Gr(k,h)$ and $\Gr(q,h)$ to \eqref{eq:Psi}.
Evaluating the resulting brackets on a $q\times h$ matrix $Z$ gives
the \emph{Gale-dual Veronese polynomial} $G_k(Z)$.
\end{definition}

The definition is independent of a chosen bracket expression: the
complementary-minor map respects the Pl\"ucker relations. Fixing volume
forms fixes its overall sign. It gives a polynomial of degree $d_0$ with
\begin{equation}\label{eq:G-character}
 G_k(gZ)=\det(g)^{k+1}G_k(Z).
\end{equation}
It vanishes on matrices of rank less than $q$.

\begin{lemma}\label{lem:G-chart}
On the chart $Z=A[I_q\mid Y]$, with $A$ invertible,
\begin{equation}\label{eq:G-chart}
 G_k(Z)=\det(A)^{k+1}\det E(Y),
\end{equation}
after fixing the sign in Definition~\ref{def:Gale}. In particular
$G_k$ is nonzero.
\end{lemma}

\begin{proof}
A basis matrix for the kernel of $[I_q\mid Y]$ is
$[-Y^{\mathsf t}\mid I_k]^{\mathsf t}=[-Y;I_k]$.
In its Veronese determinant the last $k$ rows evaluate the squares at
the coordinate points. Eliminating the square columns leaves the
matrix of the mixed monomials on the first $q$ rows, namely $E(Y)$
up to transpose and sign. Equation~\eqref{eq:G-character} gives the
factor $\det(A)^{k+1}$. If the rows of $Y$ are the vectors
$e_i+e_j$, one for each pair $i<j$, then $E(Y)$ is a permutation
matrix. This proves nonvanishing.
\end{proof}

\begin{proposition}\label{prop:equations}
For every rank-$q$ matrix $U$ satisfying $UB^{-1}U^{\mathsf t}=0$,
the polynomial $G_k(UV)$ belongs to $P$.
\end{proposition}

\begin{proof}
Complete $U$ to hyperbolic coordinates. In these coordinates write
$V=[Z^{\mathsf t}\ W^{\mathsf t}]^{\mathsf t}$ with $Z=UV$.
On the dense open subset defining the principal component,
\[
 D=V^{\mathsf t}BV=Z^{\mathsf t}W+W^{\mathsf t}Z
\]
is invertible and diagonal. If $\rank Z=q$, let $K_0$ be a basis
matrix for $\ker Z$. Then $K_0^{\mathsf t}DK_0=0$ gives a linear
relation among the $h$ vectors $\nu_2(\kappa_i)$, with nonzero
coefficients $D_{ii}$. Hence $\Psi_k(K_0)=0$, and complementary
minors give $G_k(Z)=0$. If $\rank Z<q$, the same vanishing follows
from the bracket definition. Taking closures proves the assertion.
\end{proof}

\subsection{Canonical generators of the isotropic components}

Let $X_\pm=\OGr(q,2q)^\pm$ and let $\UU$ be the tautological
rank-$q$ bundle on either component. There is a proper birational map
\begin{equation}\label{eq:isotropic-resolution}
 \pi:\Tot(\UU\otimes(\CC^h)^*)\longrightarrow Y_\pm.
\end{equation}
It is an isomorphism over the rank-$q$ locus. The rank-at-most-$(q-1)$
locus in the total space has codimension $h-q+1=k+1$; its image has
the same codimension by \eqref{eq:intersection}. Since $Y_\pm$ is
Cohen--Macaulay and smooth on the rank-$q$ locus, it is normal.
The complement of this common open set has codimension at least two
on both sides of \eqref{eq:isotropic-resolution}. Extension of
functions and of sections of reflexive canonical sheaves therefore gives
\begin{equation}\label{eq:Hartogs}
 \CC[Y_\pm]=\Gamma(\Tot(\UU\otimes(\CC^h)^*),\OO),\qquad
 \omega_{\CC[Y_\pm]}=
 \Gamma(\Tot(\UU\otimes(\CC^h)^*),\omega).
\end{equation}
No rational-singularity assertion is needed for this step.

The tangent bundle of $X_\pm$ is $\bigwedge^2\UU^*$, so
\begin{equation}\label{eq:canonical-total}
 \omega_{X_\pm}=(\det\UU)^{q-1},\qquad
 \omega_{\Tot(\UU\otimes(\CC^h)^*)}
 =\pi^*(\det\UU^*)^{k+1}.
\end{equation}
For the grading there is also a shift by the fiber dimension $qh$:
the differential of each fiber coordinate has degree one. Thus, with
$A_\pm=S/Q_\pm$ and $a=k+1$,
\begin{align}
 (A_\pm)_m&=H^0(X_\pm,\Sym^m(\UU^*\otimes\CC^h)),\label{eq:sections-ring}\\
 (\omega_{A_\pm})_{m+qh}
 &=H^0(X_\pm,\Sym^m(\UU^*\otimes\CC^h)\otimes(\det\UU^*)^a),
 \quad m\geq0.\label{eq:sections-canonical}
\end{align}
There are no canonical sections in degrees below $qh$.

For a partition $\lambda$ of $m$ of length at most $q$, the Cauchy
decomposition and the dominant-weight Borel--Weil formula give
\begin{align}
 (A_\pm)_m
 &=\bigoplus_{\lambda\vdash m}\Schur_\lambda\CC^h\otimes V_{\lambda,\pm},
 \label{eq:BW-ring}\\
 (\omega_{A_\pm})_{m+qh}
 &=\bigoplus_{\lambda\vdash m}\Schur_\lambda\CC^h
        \otimes V_{\lambda+(a^q),\pm}.\label{eq:BW-canonical}
\end{align}
Here the type-$D_q$ weights are
$(\lambda_1,\ldots,\lambda_q)$ and
$(\lambda_1,\ldots,\lambda_{q-1},-\lambda_q)$ on the two components,
with the same sign convention after the shift. The last-coordinate
convention can interchange the labels $+$ and $-$, which has no effect
on the arguments. The homogeneous-bundle formulas and the spinor
line-bundle convention are reviewed in \cite[Sections 2.2, 3.3, and 4.1]{SW}.
All the displayed weights are dominant. The formulas apply to each
Cauchy summand and do not impose an upper bound on the number $h$ of
frame columns.

\begin{lemma}\label{lem:canonical-generators}
The module $\omega_{A_\pm}$ is generated in degree $qh$, by $b_k$
elements.
\end{lemma}

\begin{proof}
Multiply sections in \eqref{eq:BW-ring} by
$H^0(X_\pm,(\det\UU^*)^a)$. On the orthogonal-group factor for a
fixed $\lambda$ this is an equivariant map to the irreducible module
$V_{\lambda+(a^q),\pm}$. It is nonzero: multiplication of a nonzero
bundle section by a nonzero scalar section on the integral variety
$X_\pm$ is nonzero. Hence it is surjective. The $\GL_h$ factors
separate the different partitions, proving generation in the initial
degree. The number of generators is
$\dim H^0(X_\pm,(\det\UU^*)^a)$. In the Weyl dimension formula for
the weight $(a,\ldots,a)$ the factors from roots $e_i-e_j$ are one,
whereas the factors from $e_i+e_j$ are
$(2a+2q-i-j)/(2q-i-j)$. This is \eqref{eq:generator-number}.
\end{proof}

Lemma~\ref{lem:linkage-canonical} now shows that $P/\II$ is generated
in degree
\[
 qh+a_R=h(k-1)=d_0
\]
and that its initial component has dimension $2b_k$.

\subsection{The explicit equations span the generators}

For fixed $V$, the function $U\mapsto G_k(UV)$ transforms by
$\det(g)^{k+1}$ under a change of row basis $U\mapsto gU$.
It therefore defines a section of $(\det\UU^*)^{k+1}$ on each
$X_\pm$. Its coefficient polynomials in $V$ induce equivariant maps
\begin{equation}\label{eq:coefficient-map}
 H^0(X_\pm,(\det\UU^*)^{k+1})^*\longrightarrow(P/\II)_{d_0}.
\end{equation}
Each source is irreducible of dimension $b_k$. The map is nonzero:
take $U=[I_q\mid0]$ and $V=[Z^{\mathsf t}\ 0]^{\mathsf t}$, where
$G_k(Z)=1$ as in Lemma~\ref{lem:G-chart}. All Gram entries of $V$
vanish, but $G_k(UV)=1$. The orthogonal symmetry interchanging the
two Grassmannian components supplies the same witness for the other
map. Thus both maps are injective.

Their images are independent modulo $\II$. Fix a maximal isotropic
column space $S_0$. Exactly one component of maximal isotropic kernels
contains planes transverse to $S_0$; every plane in the other component
meets $S_0$ nontrivially. On the latter component $\rank(U|_{S_0})<q$,
so $G_k(UV)=0$ for every frame in $S_0$. As $S_0$ ranges through its
entire Grassmannian component, this says that one coefficient module
vanishes on the corresponding $Y_\pm$. The other has a nonzero
restriction there, by the preceding witness. Restriction to the whole
$Y_\pm$ is $\SO_{2q}$-equivariant, so irreducibility makes that
restriction injective. Interchanging the two types proves independence.
The two modules consequently fill the $2b_k$-dimensional initial
component of $P/\II$, completing the proof of
Theorem~\ref{thm:ideal-intro}.

\begin{remark}
The use of the whole isotropic frame component in this argument is
essential: restriction to frames in one fixed plane is not equivariant
for the full orthogonal group. The generation statement follows from
canonical modules, rather than from a finite calculation of equations
that happen to vanish on orthogonal frames.
\end{remark}

\section{The six-by-six ideal}\label{sec:six}

For $a=(a_1,a_2,a_3)$, let
\[
 A(a)=\begin{bmatrix}
 0&-a_3&a_2\\a_3&0&-a_1\\-a_2&a_1&0
 \end{bmatrix},\qquad
 \nu(z)=(z_1^2,z_2^2,z_3^2,z_1z_2,z_1z_3,z_2z_3)^{\mathsf t}.
\]
Use the split form $B=\left[\begin{smallmatrix}0&I_3\\I_3&0\end{smallmatrix}\right]$.
For a $6\times6$ matrix $V=[v_1\ \cdots\ v_6]$, define
\begin{align}
 F_+(a,V)&=\det[\nu([I_3\mid A(a)]v_i)]_{i=1}^6,\label{eq:sixplus}\\
 F_-(a,V)&=\det[\nu([A(a)\mid I_3]v_i)]_{i=1}^6.\label{eq:sixminus}
\end{align}

\begin{theorem}\label{thm:six-explicit}
Each polynomial $F_\pm$ has degree at most eight in $a$. Write
\[
 F_\pm(a,V)=\sum_{|\alpha|\leq8}a^\alpha f_{\pm,\alpha}(V).
\]
Then $P(6,6)$ is minimally generated by the fifteen off-diagonal Gram
quadrics and the $330$ polynomials $f_{\pm,\alpha}$, each of degree twelve.
\end{theorem}

\begin{proof}
First, the polynomials vanish on the principal component. On an
orthogonal basis let $D=V^{\mathsf t}BV$; then
$VD^{-1}V^{\mathsf t}=B^{-1}$. For either isotropic matrix $U$ in
\eqref{eq:sixplus}--\eqref{eq:sixminus}, this gives
\[
 \sum_i D_{ii}^{-1}(Uv_i)(Uv_i)^{\mathsf t}=0.
\]
The six Veronese columns are therefore dependent.

To check the parameter degree, homogenize the first matrix to
$[a_0I_3\mid A(a)]$. The determinant is homogeneous of degree twelve
in $(a_0,a)$. Over $\CC(a_1,a_2,a_3)$, choose an invertible row change
that makes the third row of $A(a)$ zero. The third output coordinate
is then divisible by $a_0$, so the rows corresponding to
$z_3^2,z_1z_3,z_2z_3$ contribute a factor $a_0^4$. The determinant
of the row change on quadratic coordinates is the fourth power of
its determinant and does not involve $a_0$. Divisibility by $a_0^4$
over the fraction field implies divisibility in the polynomial ring.
Dehomogenizing proves the degree bound. The other family is obtained
by interchanging the two coordinate blocks. There are
$\binom{11}{3}=165$ possible coefficients in each family.

For independence, set
\[
 Z=\begin{bmatrix}
 1&0&0&1&1&0\\0&1&0&1&0&1\\0&0&1&0&1&1
 \end{bmatrix},\qquad
 V_+(t)=\begin{bmatrix}Z\\A(t)Z\end{bmatrix},\qquad
 V_-(t)=\begin{bmatrix}A(t)Z\\Z\end{bmatrix}.
\]
These are isotropic frames, and the Veronese determinant of $Z$ is
one. The identity
$A(a)A(t)=ta^{\mathsf t}-(a\cdot t)I_3$ gives
$\det(I_3+A(a)A(t))=(1-a\cdot t)^2$. Since
$\det\Sym^2(g)=\det(g)^4$, one obtains
\begin{align*}
 F_+(a,V_+(t))&=(1-a\cdot t)^8,&F_-(a,V_+(t))&=0,\\
 F_-(a,V_-(t))&=(1-a\cdot t)^8,&F_+(a,V_-(t))&=0.
\end{align*}
The zero entries follow from $\rank(A(a)+A(t))\leq2$.
For $|\alpha|\leq8$, the coefficient of $a^\alpha$ in the nonzero
expression is
\[
 (-1)^{|\alpha|}\frac{8!}{(8-|\alpha|)!\alpha_1!\alpha_2!\alpha_3!}
 t^\alpha.
\]
These are nonzero multiples of distinct monomials over $\CC$.
Evaluation on the two frame families proves that all $330$ classes
are independent modulo $\II(6,6)$. Theorem~\ref{thm:ideal-intro}
gives $d_0=12$ and $2b_3=330$, so they form the entire initial
generating space. This proves the assertion and minimality.
\end{proof}

\subsection{Hilbert series and the real vanishing ideal}

The canonical-module calculation is particularly concrete here.
One component of $\OGr(3,6)$ is $\PP^3$: identify the quadratic
space with $\bigwedge^2\CC^4$ and parameterize its maximal isotropic
planes by lines in $\CC^4$. If $\mathcal Q$ is the quotient bundle,
then $\UU=\OO(-1)\otimes\mathcal Q$ and $\det\UU=\OO(-2)$.
The total-space canonical bundle in \eqref{eq:canonical-total} is
the pullback of $\OO(8)$, with fiber-degree shift eighteen. This
also gives $h^0(\PP^3,\OO(8))=165$ canonical generators on each
isotropic component.

\begin{corollary}\label{cor:Hilbert}
Over $\CC$, the Hilbert series of an isotropic component ring $A_\pm$
and of the principal component ring $A=S/P(6,6)$ are
\begin{align*}
 H_{A_\pm}(t)&=\frac{1+15t+99t^2+165t^3}{(1-t)^{21}},\\
 H_A(t)&=\frac{(1+t)^{15}-2t^{12}(165+99t+15t^2+t^3)}{(1-t)^{21}}.
\end{align*}
In particular, $\deg C(6,6)=32208$.
\end{corollary}

\begin{proof}
The initial canonical degree is eighteen, so the Cohen--Macaulay ring
$A_\pm$ has $a$-invariant $-18$ and regularity $21-18=3$. Its
Hilbert numerator therefore has degree three. For completeness, on
$\PP^3$ the summands in degree $m$ have dimensions
\[
 \dim\Schur_\lambda\CC^6\cdot
 \dim\Schur_{(m,\lambda_1,\lambda_2,\lambda_3)}\CC^4,
 \qquad\lambda\vdash m,\quad\ell(\lambda)\leq3.
\]
The Weyl dimension formula in degrees zero through three gives
$1,36,645,7480$, respectively. Multiplication by $(1-t)^{21}$ gives
the first numerator. Canonical duality reverses that numerator, with
the canonical degree shift eighteen. Lemma~\ref{lem:linkage-canonical}
then places the two reversed numerators in degree twelve inside
$R$, whose Hilbert series is $(1+t)^{15}/(1-t)^{21}$. Subtracting
gives the second formula. Its numerator at one is
$2^{15}-2(165+99+15+1)=32208$.
\end{proof}

\begin{corollary}\label{cor:real}
For the positive definite real quadratic form, the real vanishing ideal
of pairwise orthogonal $h$-frames in the family \eqref{eq:triangular}
is the real form of $P(2q,h)$. In particular, Theorem~\ref{thm:six-explicit},
after a complex linear change of the ambient quadratic coordinates
and real descent, determines all polynomial equations of orthogonal
bases of $\RR^6$ and their closure.
\end{corollary}

\begin{proof}
The nonzero columns of a real orthogonal frame are independent. Since
$M\geq h$, any zero columns can be replaced by arbitrarily small
nonzero vectors in independent orthogonal directions. Thus every real
frame is a limit of the open locus defining the principal component.
Conversely, that open real locus is Zariski dense in the complex
principal component: a standard real orthogonal frame is a smooth
point, and its real smooth neighborhood has real dimension equal to
the complex dimension of the component. The analytic identity theorem
gives the required density. The complex ideal is invariant under
conjugation and therefore descends to the stated real ideal.
\end{proof}

\section{Depth, deficiency, and normality}\label{sec:depth}

We return to an arbitrary algebraically closed field of characteristic
different from two, retaining \eqref{eq:triangular} and
\eqref{eq:homological-notation}. The arguments in this section do not
use the characteristic-zero representation decomposition.

\subsection{The deficiency module}

\begin{proposition}\label{prop:Ext}
For $A=S/P$, the only nonzero Ext modules into $S(-L)$ are
\begin{equation}\label{eq:Ext}
 \Ext^e_S(A,S(-L))=(J/\II)(a_R),\qquad
 \Ext^{e+1}_S(A,S(-L))=Z_0(a_R).
\end{equation}
\end{proposition}

\begin{proof}
Set $T=P/\II$. Lemma~\ref{lem:linkage-canonical} says that $T$ and
$R=S/\II$ are Cohen--Macaulay modules of dimension $N$. Apply
$\Ext_S^\bullet(-,S(-L))$ to
\[
 0\longrightarrow T\longrightarrow R\longrightarrow A\longrightarrow0.
\]
Canonical duality on the two Cohen--Macaulay component rings gives
\begin{equation}\label{eq:dual-sequence}
 0\longrightarrow\omega_A\longrightarrow R(a_R)
 \longrightarrow(A_+\oplus A_-)(a_R)
 \longrightarrow\Ext^{e+1}_S(A,S(-L))\longrightarrow0,
\end{equation}
and all other Ext modules vanish.

After removing the common shift, each component of the middle map
is a graded $R$-linear map $R\to A_\pm$. Its value on one is a
scalar. The scalar is nonzero, since at the generic point of $Y_\pm$
the inclusion $T\to R$ is an isomorphism. Normalize the two canonical
identifications to make both scalars one. The middle map is then the
pair of quotient maps. Its kernel is $J/\II$ and its cokernel is
$Z_0$, by \eqref{eq:union}. This proves \eqref{eq:Ext}.
\end{proof}

\begin{corollary}\label{cor:local-depth}
The locus $Z=\Spec Z_0$ is contained in $C$. The ring $A$ is
Cohen--Macaulay away from $Z$, and at every $x\in Z$ it has local
depth $\dim\OO_{C,x}-1$. It satisfies $S_k$ but not $S_{k+1}$.
\end{corollary}

\begin{proof}
The support of an Ext module is contained in the support of its first
argument. The second module in \eqref{eq:Ext} therefore proves
$Z\subset C$. Localize in the regular ring $S$ at a prime lying
on $C$. If it is not on $Z$, the projective dimension of $A$ there
is $e$; if it is on $Z$, the nonzero last Ext module makes that
projective dimension $e+1$. The codimension of the principal prime
remains $e$ under this localization. Auslander--Buchsbaum gives the
two asserted depths.

The locus $Z$ has codimension $k+1$ by \eqref{eq:intersection}. Thus
at points of $Z$ the local dimension is at least $k+1$, so the depth
is at least $k$. Elsewhere the ring is Cohen--Macaulay. This proves
$S_k$. At the generic point of $Z$, local dimension is $k+1$ and
depth is $k$, proving failure of $S_{k+1}$.
\end{proof}

\begin{proposition}\label{prop:graded-endpoints}
The global invariants and last free module are
\[
 \depth A=N-1,\qquad\pd_S A=e+1,\qquad\reg A=e-1,
 \qquad F_{e+1}=S(-2e).
\]
Moreover, $\depth\omega_A=N-k+1$.
\end{proposition}

\begin{proof}
The first two assertions follow from Proposition~\ref{prop:Ext}.
The module $Z_0(a_R)$ is cyclic, with its generator in degree $-a_R$.
Dualizing a minimal graded free resolution preserves minimality of
the differentials. Thus the last free module has rank one. If its
shift is $j$, its dual into $S(-L)$ has generator degree $L-j$.
Equation~\eqref{eq:Ext} gives $L-j=-a_R$, hence $j=2e$.

The initial degree of $J/\II$ is two. Indeed, the diagonal norm
quadrics belong to $J$ and are independent of the off-diagonal
quadrics by column multidegrees. There are no linear equations on
the isotropic union: in split coordinates every coordinate matrix
unit is an isotropic frame. Graded local duality therefore gives
\[
 \operatorname{end} H^{N-1}_{\mathfrak m}(A)=a_R,
 \qquad
 \operatorname{end} H^N_{\mathfrak m}(A)=a_R-2,
\]
and all other local cohomology modules vanish. Since $a_R+N=e$,
their degree formula for regularity gives $\reg A=e-1$.

From \eqref{eq:union} one has $\depth(S/J)=N-k$.
The sequence $0\to J/\II\to R\to S/J\to0$ then gives
$\depth(J/\II)=N-k+1$. Apply the first isomorphism in
\eqref{eq:Ext} to obtain the last assertion.
\end{proof}

For $(M,h)=(6,6)$, these formulas give dimension $21$, depth $20$,
regularity $14$, projective dimension $16$, and last free module
$S(-30)$. Its canonical module has depth $19$.

\subsection{The boundary determinant}

Corollary~\ref{cor:local-depth} gives $S_2$. To prove normality it
remains to prove regularity in codimension one.

\begin{lemma}\label{lem:det-irreducible}
Let $q=\binom k2$, $k\geq3$, and set $D_k(Y)=\det E(Y)$ for a
$q\times k$ matrix $Y$. Then $D_k$ is irreducible and its generic
zero has $\rank E(Y)=q-1$. At a general point of $D_k=0$, the
relation $E(Y)\lambda=0$ has all $\lambda_r\ne0$, and
\[
 Y^{\mathsf t}\diag(\lambda)Y=\diag(\mu_1,\ldots,\mu_k)
 \quad\text{has all }\mu_j\ne0.
\]
\end{lemma}

\begin{proof}
Let $H$ be the $q$-dimensional space spanned by the hollow quadrics
$x_i x_j$, $i<j$. Define the incidence scheme
\[
 \mathcal Z=
 \{([Q],y_1,\ldots,y_q)\in\PP(H)\times(\mathbb A^k)^q:
                         Q(y_r)=0\text{ for all }r\}.
\]
It is the $q$-fold fiber product of the universal affine quadric
over $\PP(H)$. To check flatness, put
$W_m=K[x_1,\ldots,x_k]_m$, with $W_m=0$ for $m<0$.
Its degree-$m$ coordinate module is the cokernel of multiplication
by the universal quadric,
\[
 \OO_{\PP(H)}(-1)\otimes_K W_{m-2}
 \longrightarrow \OO_{\PP(H)}\otimes_K W_m.
\]
Every fiber map is injective, since its defining quadratic is nonzero.
The map has constant rank, so its cokernel is locally free. Taking
the direct sum over $m$ proves flatness of the universal affine
quadric, and hence of $\mathcal Z$ over $\PP(H)$.

The hollow quadric $x_1x_2+x_1x_3+x_2x_3$ has rank three in every
characteristic different from two. Thus a general hollow quadric
has rank at least three and is geometrically integral, as is its
$q$-fold product. Flatness over the integral base makes the generic
fiber dense and makes each affine coordinate ring inject into its
generic-fiber algebra, which is a domain. It follows that
$\mathcal Z$ is integral.

Projection to $Y$ is proper because the other factor is projective.
A point is in its image precisely when there is a nonzero hollow
quadric vanishing at all its rows, or equivalently when
$\ker E(Y)^{\mathsf t}\ne0$. The image is therefore exactly $D_k=0$,
which is irreducible as a closed set. The choice $y_r=e_i+e_j$,
one row for each pair, makes $E(Y)$ the identity. Thus $D_k$ is
nonzero and is a scalar times a power of an irreducible polynomial.

We show that this power is one. Fix a general integral quadric $Q$,
with coefficient vector $a\in H$. Evaluations of $H$ on $Q$ span
$a^\perp$: the only degree-two polynomial in $H$ vanishing on $Q$
is a scalar multiple of $Q$. Choose $q$ smooth points of $Q$ such
that every $q-1$ of their evaluation vectors are independent.
Such tuples form a nonempty open subset of the irreducible product
$V(Q)^q$, since each of the finitely many independence conditions
holds on a nonempty open subset. At this tuple, $E$ has rank $q-1$.
Its left kernel is spanned by $a$, and its right kernel is spanned
by a vector $\lambda$ with every $\lambda_r\ne0$. Consequently
\[
 \operatorname{adj}(E)=\gamma\lambda a^{\mathsf t}
 \qquad\text{for some }\gamma\ne0.
\]
If only the row $y_r$ moves, in the direction $v$, differentiation gives
\[
 dD_k(Y)(0,\ldots,v,\ldots,0)
       =\gamma\lambda_r\,dQ_{y_r}(v).
\]
Choose $v$ transverse to $Q$ at $y_r$. The right side is nonzero,
so $D_k$ is reduced at this point and cannot be a proper power.
This proves irreducibility as a polynomial and also the generic
rank and full-support assertions.

For each $j$, the restrictions of $H+\langle x_j^2\rangle$ to $Q$
have dimension $q$: their only relation is $Q$, and $x_j^2$ is
not hollow. Evaluation at $q$ general points is an isomorphism
on this restriction space. The relation $\lambda$ among the hollow
evaluations therefore does not annihilate $x_j^2$, so
\[
 \mu_j=\sum_r\lambda_r y_{rj}^2\ne0.
\]
These finitely many evaluation conditions can be imposed together
with the preceding independence conditions. They give a nonempty
open subset of $D_k=0$ on which all the stated conclusions hold.
\end{proof}

\begin{proposition}\label{prop:boundary-R1}
The principal component is smooth at the generic point of every
divisor it has in the rank-$q$ isotropic stratum.
\end{proposition}

\begin{proof}
Use Lemma~\ref{lem:chart} with $m=0$. The equations are
$E(Y)\lambda=0$ and the isotropic component is $\lambda=0$.
At a general point of $D_k=0$, invert a $(q-1)$-minor and eliminate
$q-1$ of the $\lambda$ coordinates. The remaining equation is
\[
 D_k(Y)u=0
\]
up to a unit. The branch $u=0$ is isotropic. On the other branch,
Lemma~\ref{lem:det-irreducible} gives points with $u\ne0$ for which
all first $q$ norms $\lambda_r$ and all remaining norms $-\mu_j$
are nonzero. Thus the branch $D_k=0$ is the principal component.
Where $dD_k\ne0$, it is smooth also at $u=0$, its intersection with
the isotropic branch.

For completeness, this accounts for every divisor in the indicated
stratum. If $D_k$ is a unit, the chart has only the isotropic
component, so no point there belongs to $C$. On the isotropic
stratum itself, where $\lambda=0$, the possible intersection with
$C$ is therefore supported on $D_k=0$. The preceding calculation
places the generic point of this irreducible hypersurface in $C$;
closedness places its whole closure there. The locus where the
rank, nonvanishing, or smoothness conditions used above fail is a
proper closed subset of that hypersurface. Since the hypersurface
has dimension $N-1$ after adjoining the free chart coordinates,
the exceptional subset has dimension at most $N-2$. The column
and Witt charts cover the rank-$q$ isotropic stratum, proving the
assertion.
\end{proof}

\subsection{The remaining codimension-one strata}

For a frame let $t$ be the number of columns of nonzero norm and $r$
the rank of its isotropic columns. For a fixed choice of the $t$
indices, the stratum has dimension
\begin{equation}\label{eq:strata}
 D(t,r)=Mt-\frac{t(t-1)}2
       +r(M+h-2t-r)-\frac{r(r+1)}2,
 \quad t+r\leq h,\quad t+2r\leq M.
\end{equation}
The first term chooses the nonisotropic orthogonal columns. In their
nondegenerate orthogonal complement, the orthogonal Grassmannian of
isotropic $r$-planes has dimension $r(M-t-r)-r(r+1)/2$.
A spanning $(h-t)$-tuple in the chosen plane contributes $r(h-t)$
parameters. These choices give \eqref{eq:strata}; every component
of the stratum has this dimension.

Suppose $k\geq4$, put $x=q-r$, and write $c_x(t)=N-D(t,q-x)$.
At our parameters,
\begin{equation}\label{eq:strata-codimension}
 c_x(t)=\frac{2kx+t^2-4tx-t+3x^2-x}{2}.
\end{equation}
The constraints become $0\leq t\leq\min(2x,k+x)$, and
\[
 c_x(t+1)-c_x(t)=t-2x.
\]
Thus the codimension decreases through the allowed successive values.
If $1\leq x\leq k$, set $t=2x-s$. Then
\[
 c_x(t)=\frac{x(2k-x-3)}2+\frac{s(s+1)}2.
\]
The first term is concave in $x$; its endpoint values are $k-2$
and $k(k-3)/2$, both at least two. None of these strata can
contribute a divisor.

If $x\geq k$, set $t=k+x-s$ and $r=q-x$. Now
\[
 c_x(t)=r+\frac{s(2x-2k+s+1)}2.
\]
For $r\geq2$ this is at least two. For $r=1$, the value $s=0$
gives the codimension-one stratum $t=h-1$, whose frames have
rank $t+r=h$; if $s\geq1$, the codimension is at least
$1+q-k\geq3$. For $r=0$, the value $s=0$ is the full-rank open
stratum, whereas $s\geq1$ gives codimension at least $q-k+1\geq3$.
Finally, $x=0$ forces $t=0$ and $r=q$, the maximal-isotropic
stratum treated in Proposition~\ref{prop:boundary-R1}.

A full-column-rank frame is smooth on $\XX$: a relation among the
Jacobian rows is a hollow symmetric matrix $K_0$ with $VK_0=0$,
which forces $K_0=0$. This treats all remaining divisors for $k\geq4$.

For $k=3$, one has $q=3$, $h=6$, and $N=21$. Substitution in
\eqref{eq:strata} gives the complete list of strata with
$N-D(t,r)\leq1$:
\begin{center}
\begin{tabular}{@{}c c c c@{}}
\toprule
$t$&$r$&$N-D(t,r)$&$\rank V=t+r$\\
\midrule
0&3&0&3\\
2&2&1&4\\
4&1&1&5\\
5&0&1&5\\
6&0&0&6\\
\bottomrule
\end{tabular}
\end{center}
The first and last rows have already been treated. In the
$(4,1)$ and $(5,0)$ cases, $\ker V$ is one-dimensional.
Every symmetric Jacobian relation has image in this kernel and
therefore has the form $b uu^{\mathsf t}$ for a generator $u$.
Some coordinate $u_i$ is nonzero, and the zero diagonal gives
$b u_i^2=0$, hence $b=0$.

For $(t,r)=(2,2)$, split off the two nonisotropic columns.
Write the remaining four columns in a basis of their isotropic
two-plane. For an explicit rank-two coordinate matrix use
\[
 [I_2\mid Y],\qquad
 Y=\begin{bmatrix}1&1\\1&-1\end{bmatrix}.
\]
The first two coordinates of any relation among all six columns
are zero, because the nonisotropic span meets its orthogonal
complement trivially. The residual kernel has basis matrix
$W=[-Y;I_2]$, with rows
\[
 (-1,-1),\quad(-1,1),\quad(1,0),\quad(0,1).
\]
A symmetric relation on the residual four columns has the form
$WTW^{\mathsf t}$, where
$T=\left[\begin{smallmatrix}a&b\\b&c\end{smallmatrix}\right]$.
Its last two diagonal entries force $a=c=0$, and its first
diagonal entry then forces $2b=0$. Since the characteristic
is different from two, $T=0$, so there is no nonzero Jacobian
relation at this frame.

This coordinate matrix can be placed in every isotropic two-plane.
For each such plane, the spanning four-tuples form an irreducible
open subset of $\Mat_{2,4}$, and full Jacobian rank is a nonempty
open condition there. The argument applies to both isotropic-plane
types in each fixed orthogonal complement. The bad locus has
positive codimension in every family of spanning four-tuples, so
it has dimension at most $N-2$ in the stratum. Thus it contributes
no divisor of $C$.

We have proved $R_1$ in all cases. Together with $S_2$ from
Corollary~\ref{cor:local-depth}, Serre's criterion proves normality.
Propositions~\ref{prop:Ext} and \ref{prop:graded-endpoints} now complete
the proof of Theorem~\ref{thm:depth-intro}.

\section{Generic transverse singularities in odd ambient dimension}\label{sec:transverse}

We now use a different parameter range. Let
\begin{equation}\label{eq:odd-parameters}
 k\geq2,\qquad 2\leq c\leq k,\qquad
 h=\binom k2+c,\qquad q=h-k,\qquad M=2q+1.
\end{equation}
Formula~\eqref{eq:prime-threshold} gives $M=\Dpr(h)$.
Let $S_{\mathrm{iso}}$ be the stratum of frames spanning a maximal
isotropic $q$-plane. It is irreducible, and its dimension and codimension
in $\XX(M,h)$ are
\begin{equation}\label{eq:odd-stratum}
 \dim S_{\mathrm{iso}}=qh+\frac{q(q+1)}2,
 \qquad\codim_{\XX(M,h)}S_{\mathrm{iso}}=c.
\end{equation}
The latter follows by subtracting from $Mh-\binom h2$.

\begin{theorem}\label{thm:transverse}
Let $L$ be the function field of $S_{\mathrm{iso}}$, and set
\[
 \delta=k-c,\qquad a=2\delta-k=k-2c=k^2-2h.
\]
At the generic point of $S_{\mathrm{iso}}$, the local ring of
$\XX(M,h)$ is
\begin{equation}\label{eq:transverse-ring}
 \left(L[z_1,\ldots,z_k]/(Q_1,\ldots,Q_\delta)\right)_{(z)},
\end{equation}
where the $Q_i$ are hollow quadrics and their projective intersection
is geometrically smooth. If $\delta=0$, this local ring is regular.
For $\delta>0$ the classification is
\begin{center}
\begin{tabular}{@{}lll@{}}
\toprule
Condition & Characteristic zero & Odd characteristic\\
\midrule
$a<0$ & rational and canonical & strongly $F$-regular\\
$a=0$ & log canonical, not rational & $F$-pure, not $F$-rational\\
$a>0$ & not log canonical & not $F$-pure\\
\bottomrule
\end{tabular}
\end{center}
Every assertion in this table concerns the generic transverse local ring.
\end{theorem}

\subsection{Elimination and dominance}

Apply Lemma~\ref{lem:chart} with $m=1$ and write the last row of $z$
as $(z_1,\ldots,z_k)$. The residual equations are
\begin{equation}\label{eq:odd-residual}
 z_i z_j-\sum_{r=1}^q\lambda_r y_{ri}y_{rj}=0,
 \qquad i<j.
\end{equation}
Put $e'=\binom k2$, so $e'-q=\delta$. The matrix $E(Y)$ has
generic rank $q$: choose its $q$ rows of $Y$ among distinct vectors
$e_i+e_j$, so its columns are distinct coordinate vectors. At the
generic point of the stratum, $q$ equations eliminate the $\lambda_r$.
The remaining quadrics form the kernel of $E(Y)^{\mathsf t}$ in the
$e'$-dimensional space
\[
 H=\langle z_i z_j:i<j\rangle.
\]
All chart parameters other than $\lambda,z$ become elements of $L$.
Consequently this is the exact local presentation
\eqref{eq:transverse-ring}, not only an associated graded ring.

\begin{lemma}\label{lem:dominance}
If $\delta>0$, the rational map
\[
 \Mat_{q,k}\dashrightarrow\Gr(\delta,H),\qquad
 Y\longmapsto\ker E(Y)^{\mathsf t}
\]
is dominant. A general $\delta$-dimensional subspace of $H$ cuts
out a geometrically smooth projective complete intersection of
dimension $c-1$.
\end{lemma}

\begin{proof}
We first show that the smooth locus in the indicated Grassmannian
is nonempty. Choose $\delta$ diagonal quadrics
$\sum_i b_{ji}z_i^2$ such that every maximal minor of the
$\delta\times k$ matrix $(b_{ji})$ is nonzero. Vandermonde columns
over the infinite field supply such a matrix. Any common projective
zero has at least $\delta+1$ nonzero coordinates; otherwise a
full-column-rank submatrix of $(b_{ji})$ would annihilate the nonzero
vector of their squares. The Jacobian therefore has rank $\delta$
at every common zero, because two is invertible. The intersection
is a nonempty smooth complete intersection of dimension
$k-1-\delta=c-1\geq1$.

The Koszul resolution shows that it is connected and has no linear
equations. Smoothness and connectedness make it integral. It contains
$k$ linearly independent projective points. Send these points to the
coordinate points by a projective change of coordinates. The resulting
quadrics are hollow and still define a smooth complete intersection.
Smoothness of such an intersection is an open condition on the
Grassmannian, proving the required nonempty open subset.

Now let $K_0\subset H$ lie in that open subset, and let $T$ be its
projective zero set. Its degree-two ideal is exactly $K_0$, since
it is a positive-dimensional complete intersection of those quadrics.
It follows that the hollow evaluation vectors
$(z_i z_j)_{i<j}$ at points of $T$ span $K_0^\perp$, a vector space
of dimension $q$. Choose $q$ points with independent evaluations
and use representatives as the rows of $Y$. Then $E(Y)$ has rank
$q$ and $\ker E(Y)^{\mathsf t}=K_0$. The image contains the whole
nonempty open subset, proving dominance. These arguments also apply
after algebraically closed extensions, giving geometric smoothness.
\end{proof}

The quadrics in \eqref{eq:transverse-ring} are consequently a regular
sequence, and their graded ring is a normal complete intersection
of dimension $c\geq2$, singular only possibly at the vertex. Its
$a$-invariant is $2\delta-k=a$.

\subsection{Characteristic zero}

Suppose $\delta>0$, and let $T\subset\PP^{k-1}_L$ be the smooth
intersection from Lemma~\ref{lem:dominance}. Its cone is resolved by
blowing up the vertex; the resolution is $\Tot(\OO_T(-1))$.
Adjunction gives $\omega_T=\OO_T(a)$. If $E$ is the exceptional
divisor, then $\OO_E(E)=\OO_T(-1)$, and the discrepancy is
\begin{equation}\label{eq:discrepancy}
 -a-1.
\end{equation}
For example, this follows by restricting
$(K_{\Tot(\OO_T(-1))}+E)|_E=K_T$ and using that the cone is
Gorenstein. The resolution has one smooth exceptional divisor.
Thus the cone is log canonical exactly when $a\leq0$, and canonical
exactly when $a<0$.

Rationality can be read off directly. Since $T$ is an arithmetically
Cohen--Macaulay complete intersection, its intermediate cohomology
$H^i(T,\OO_T(m))$ vanishes for $m\geq0$ and
$0<i<\dim T$. By Serre duality its top cohomology is
\[
 H^{c-1}(T,\OO_T(m))
   =H^0(T,\OO_T(a-m))^*.
\]
All these groups vanish for $m\geq0$ exactly when $a<0$. If
$a\geq0$, the group at $m=0$ is nonzero. The cohomology of the
total-space resolution is the direct sum of these groups over
$m\geq0$, which computes the higher direct images because the
cone is affine. This proves rationality precisely for $a<0$ and
establishes the characteristic-zero column of the table.

\subsection{Odd characteristic}

We recall the complete-intersection Frobenius criteria in the form
needed here. Let
$R=L[z]_{(z)}/(Q_1,\ldots,Q_\delta)$, where $L$ is $F$-finite,
and put $F=Q_1\cdots Q_\delta$. Fedder's criterion \cite{Fedder}
says that $R$ is $F$-pure if and only if
\begin{equation}\label{eq:Fedder}
 F^{p-1}\notin(z_1^p,\ldots,z_k^p).
\end{equation}
If $u$ is a nonzero element for which $R_u$ is strongly $F$-regular,
Glassbrenner's criterion \cite[Theorem 4.1]{Glassbrenner} gives strong
$F$-regularity from
\begin{equation}\label{eq:Glassbrenner}
 uF^{p^e-1}\notin(z_1^{p^e},\ldots,z_k^{p^e})
 \quad\text{for some }e\geq1.
\end{equation}
Conversely, in a strongly $F$-regular ring this nonmembership occurs
for every nonzero $u$ and a suitable $e$. These forms follow from
the usual complete-intersection colon-ideal formula. Our residue
field $L$ is finitely generated over the perfect field $K$, hence
is $F$-finite. The cone is regular away from its vertex, so
localization at a nonzero linear coordinate is regular.

Suppose $2\delta<k$. Work first with arbitrary independent hollow
quadrics as parameters. At the special system
\[
 Q_i=z_{2i}z_{2i+1},\qquad 1\leq i\leq\delta,
\]
the coefficient of
$z_1\prod_{j=2}^{2\delta+1}z_j^{p-1}$ in $z_1F^{p-1}$ is one.
It is a polynomial in the coefficients of the $Q_i$, so it is
generically nonzero. Every exponent is below $p$. The nonvanishing
locus meets the open locus of smooth projective complete intersections
from Lemma~\ref{lem:dominance}. On this intersection, apply
\eqref{eq:Glassbrenner} with $u=z_1$. Dominance then proves strong
$F$-regularity of the generic transverse ring. The monomial system
is only a witness that a coefficient polynomial is nonzero; no
smoothness assertion about that special system is used.

If $2\delta=k$, use the special system
$Q_i=z_{2i-1}z_{2i}$. The coefficient of
$\prod_j z_j^{p-1}$ in $F^{p-1}$ is one at that system and hence
nonzero generically. This proves $F$-purity by \eqref{eq:Fedder}.
For every $p$-power $Q$ and nonzero linear coordinate $u$,
\[
 \deg(uF^{Q-1})=1+k(Q-1)>k(Q-1).
\]
Every monomial therefore lies in $(z_1^Q,\ldots,z_k^Q)$, contradicting
the necessary direction of \eqref{eq:Glassbrenner}. The ring is not
strongly $F$-regular. An $F$-rational $F$-finite Gorenstein local
ring is strongly $F$-regular \cite[Corollary 4.7 and Theorem 5.5]{HH};
thus it is not $F$-rational.

Finally, if $2\delta>k$, the same degree comparison already places
$F^{p-1}$ in $(z_1^p,\ldots,z_k^p)$. It is not $F$-pure, and again
cannot be $F$-rational. This proves the remaining column of
Theorem~\ref{thm:transverse}. Changes of basis among the quadrics do
not affect these criteria, so the nonempty open conditions used above
descend from the space of ordered systems to the Grassmannian.

\section{Density, factorial examples, and the questions answered}\label{sec:density}

\subsection{The counterexample interval}

Put $T_k=\binom k2$. A direct simplification of
\eqref{eq:prime-threshold} gives
\begin{equation}\label{eq:prime-blocks}
 \begin{aligned}
 \Dpr(h)&=2h-2k+1&&\text{if }T_k+2\leq h\leq T_{k+1},\\
 \Dpr(T_{k+1}+1)&=2(T_{k+1}+1)-2k.
 \end{aligned}
\end{equation}
These identities follow by bounding the square roots between consecutive
integers; in the first interval
$(2k-1)^2<8h-7<(2k+1)^2$, with the endpoints of the integer floors
giving the displayed value.

For $T_k+2\leq h\leq\lfloor k^2/2\rfloor$, write $c=h-T_k$.
Then $2\leq c\leq\lfloor k/2\rfloor$, so
Theorem~\ref{thm:transverse} applies and its invariant is
\[
 a=k-2c=k^2-2h\geq0.
\]
The variety at $M=\Dpr(h)$ is a normal integral complete intersection
by \cite[Theorems 1 and 2]{CS}, but its indicated local ring is not
rational in characteristic zero and not $F$-rational in odd
characteristic. For strict inequality $h<k^2/2$, the invariant is
positive and gives failure of log canonicity and of $F$-purity.

The interval contains $\lfloor k/2\rfloor-1$ integers. Through the
completed block ending at $T_{K+1}+1$, their total number is
\[
 \sum_{k=4}^K\left(\left\lfloor\frac k2\right\rfloor-1\right)
 =\left\lfloor\frac{(K-2)^2}{4}\right\rfloor.
\]
The endpoint is $K^2/2+O(K)$; the ratio tends to $1/2$. An incomplete
last block changes the count by $O(K)$, proving natural density.
The equality cases $a=0$ occur only once for each even $k$, so
removing them leaves density $1/2$ for the strict failures as well.

\subsection{Factoriality}

The UFD bound in \cite[Theorem 3]{CS}, for $h\geq4$, is the minimum
of an even candidate and the odd candidate
\begin{equation}\label{eq:UFD-odd}
 2\left\lceil\frac{2h-1-\sqrt{8h-31}}2\right\rceil+1.
\end{equation}
For $4\leq c\leq\lfloor k/2\rfloor$,
\[
 8h-31=(2k-1)^2+8c-32,
 \qquad 2k-1\leq\sqrt{8h-31}<2k+1.
\]
Expression~\eqref{eq:UFD-odd} is therefore $2h-2k+1=M$, so the
frame ring is factorial. Removing $c=2,3$ from each block removes
only $O(\sqrt H)$ integers up to $H$. The factorial subfamily, and
its strict non-log-canonical and non-$F$-pure subfamily, still have
natural density $1/2$. This completes the proof of
Theorem~\ref{thm:negative-intro}.

\subsection{Examples and scope}

When $a=0$, write $k=2c$. Then
\begin{equation}\label{eq:CY-boundary}
 h=2c^2,\qquad M=(2c-1)^2,\qquad\delta=c.
\end{equation}
The generic transverse projective variety is a smooth intersection
of $c$ quadrics in $\PP^{2c-1}$ with trivial canonical bundle.
For $c=2$ it is an elliptic quartic curve; for $c=3$ it is a K3
surface of degree eight; for $c=4$ it is a Calabi--Yau threefold.
For $c\geq4$, the ambient frame ring is factorial. In odd
characteristic these generic transverse cones are $F$-pure but
not $F$-rational.

\begin{center}
\begin{tabular}{@{}r r c l@{}}
\toprule
$M$&$h$&$(k,c)$&Characteristic-zero conclusion\\
\midrule
9&8&$(4,2)$&not rational; elliptic transverse cone\\
15&12&$(5,2)$&not log canonical\\
25&18&$(6,3)$&not rational; K3 transverse cone\\
49&32&$(8,4)$&factorial, not rational\\
63&40&$(9,4)$&factorial, not log canonical\\
\bottomrule
\end{tabular}
\end{center}

The two families answer distinct questions. At the triangular
complete-intersection parameters, the principal component has an
explicit defining ideal and is normal with a controlled failure of
Cohen--Macaulayness. At the prime parameters in
\eqref{eq:bad-range}, the full frame variety is already normal and
integral, yet its generic transverse cones obstruct rationality.
Thus neither normality nor factoriality supplies the proposed
rational-singularity conclusion.

Several classification problems remain. The normality of all principal
components is not decided by the triangular family, and their defining
ideals outside that family remain to be determined. Likewise,
$a<0$ in Theorem~\ref{thm:transverse} controls only the generic
maximal-isotropic stratum; it gives no assertion about every special
point or about other strata. Determining the exact rational-singularity
threshold requires this additional local analysis.

\paragraph{Acknowledgment.}
The author acknowledges OpenAI GPT-6 for assistance with research
exploration, proof development, verification code, and manuscript
preparation.

\begin{thebibliography}{8}

\bibitem{BH}
W.~Bruns and J.~Herzog,
\emph{Cohen--Macaulay Rings}, revised ed.,
Cambridge Studies in Advanced Mathematics, vol.~39,
Cambridge University Press, Cambridge, 1998.
\doi{10.1017/CBO9780511608681}.

\bibitem{CS}
L.~Casabella and A.~Sammartano,
\emph{The variety of orthogonal frames},
preprint, 2025.
\href{https://arxiv.org/abs/2512.25058v1}{arXiv:2512.25058v1}.

\bibitem{Fedder}
R.~Fedder,
\emph{$F$-purity and rational singularity},
Trans. Amer. Math. Soc. \textbf{278} (1983), 461--480.
\doi{10.1090/S0002-9947-1983-0701505-0}.

\bibitem{Glassbrenner}
D.~Glassbrenner,
\emph{Strong $F$-regularity in images of regular rings},
Proc. Amer. Math. Soc. \textbf{124} (1996), 345--353.
\doi{10.1090/S0002-9939-96-03030-4}.

\bibitem{HH}
M.~Hochster and C.~Huneke,
\emph{$F$-regularity, test elements, and smooth base change},
Trans. Amer. Math. Soc. \textbf{346} (1994), 1--62.
\doi{10.1090/S0002-9947-1994-1273534-X}.

\bibitem{HJPS}
M.~Hochster, J.~Jeffries, V.~Pandey, and A.~K.~Singh,
\emph{When are the natural embeddings of classical invariant rings pure?},
Forum Math. Sigma \textbf{11} (2023), e67, 43~pp.
\doi{10.1017/fms.2023.67}.
\href{https://arxiv.org/abs/2210.09351v3}{arXiv:2210.09351v3}.

\bibitem{PS}
C.~Peskine and L.~Szpiro,
\emph{Liaison des vari\'et\'es alg\'ebriques. I},
Invent. Math. \textbf{26} (1974), 271--302.
\doi{10.1007/BF01425554}.

\bibitem{SW}
S.~V.~Sam and J.~Weyman,
\emph{Littlewood complexes and analogues of determinantal varieties},
Int. Math. Res. Not. IMRN (2015), 4663--4707.
\doi{10.1093/imrn/rnu078}.
\href{https://arxiv.org/abs/1303.0546v3}{arXiv:1303.0546v3}.

\end{thebibliography}
\end{document}
